\subsection{One clamped--free magnetic beam throughout the computational study}
\label{sec:unified-model}
We instantiate the rod reduction as a straight, planar, small-slope
Euler--Bernoulli cantilever with a concentrated magnetic dipole at its free
tip. The base position and slope are fixed. The elastic part is linear in
curvature; the full trigonometric magnetic potential is retained. These
choices define a reproducible mechanical specialization of the spatial
rod, rather than a catheter-specific material identification. Table~\ref{tab:unified-si}
states all dimensional parameters. Gravity, imposed fluid flow and field
gradient are zero in this specialization. Still-fluid/structural losses
are represented by a mass-proportional Rayleigh term. The more general
force accounting above remains applicable when those terms are present.

\begin{table}[t]
\centering
\caption{Declared parameters of the unified magnetic cantilever. These are
illustrative mechanical parameters, not measurements from a clinical or
spaceflight device.}
\label{tab:unified-si}
\setlength{\tabcolsep}{3pt}
\begin{tabular}{lll}
\toprule
Quantity & Symbol & Value\\
\midrule
Length & $L$ & $0.08\,\mathrm m$\\
Circular radius & $r_c$ & $0.0006\,\mathrm m$\\
Young's modulus & $E$ & $5.0\times10^6\,\mathrm{Pa}$\\
Density & $\rho$ & $1100\,\mathrm{kg\,m^{-3}}$\\
Tip dipole & $\mu$ & $2.0\times10^{-5}\,\mathrm{A\,m^2}$\\
Axial bias field & $B_0$ & $0.03\,\mathrm T$\\
Mass damping rate & $\alpha$ & $0.55558524\,\mathrm{s^{-1}}$\\
Transverse field amplitude & $B_a$ & $0.006\,\mathrm T$\\
Drive frequency & $\Omega_d$ & $9.4449490\,\mathrm{rad\,s^{-1}}$\\
Ramp time & $t_r$ & $0.25\,\mathrm s$\\
Uniform Hermite elements & $N_e$ & $32$\\
\bottomrule
\end{tabular}
\end{table}

Write $w(s,t)$ for transverse displacement and
$\vartheta(t)=w_s(L,t)$ for the small-slope tip rotation. A uniform field
$(B_0,B_y(t))$ has zero translational dipole force and the tip moment
\begin{equation}
 M_{\rm tip}=\mu\{B_y(t)\cos\vartheta-B_0\sin\vartheta\}.
 \label{eq:unified-tip-torque}
\end{equation}
This is the natural moment boundary load. It does not impose a second
kinematic constraint at the free tip. The magnetic potential is
$-\mu(B_0\cos\vartheta+B_y\sin\vartheta)$, so that
\eqref{eq:unified-tip-torque} follows by variation. Keeping this potential
exact does not remove the small-slope approximation in the beam kinematics.

Let $q=(w_1,\vartheta_1,\ldots,w_{N_e},\vartheta_{N_e})^{\mathsf T}$
contain the unconstrained cubic-Hermite nodal coordinates, $e_\vartheta$
select its final rotation, and $M,K$ be the consistent mass and bending
stiffness matrices. The physical equations and static-bias energy are
\begin{align}
 M\ddot q+\alpha M\dot q+Kq
 &=e_\vartheta M_{\rm tip},
 \label{eq:unified-fe}\\
 H_0(q,p)&=\tfrac12p^{\mathsf T}M^{-1}p+\tfrac12q^{\mathsf T}Kq\notag\\
 &\quad+E_0\{1-\cos(e_\vartheta^{\mathsf T}q)\},
 \label{eq:unified-SI-H}\\
 E_0&=\mu B_0=6\times10^{-7}\,\mathrm J,\qquad p=M\dot q.
\end{align}
There is no axial end restraint and no artificial stretching quartic.
The nonlinearity used in learning and in the KAM calculation is the
magnetic nonlinearity of this same cantilever.

The static-field linearization is
$K_b=K+E_0e_\vartheta e_\vartheta^{\mathsf T}$.
Choose $\Phi^{\mathsf T}M\Phi=I$ and
$\Phi^{\mathsf T}K_b\Phi=\operatorname{diag}(\omega_j^2)$,
with positive tip-slope entries fixing modal signs. For $n$ retained modes,
$q=\Phi_n a$, and the modal momentum is $\pi=\dot a$.
Set $\omega_0=\omega_1$, $\tau=\omega_0t$, $\nu_j=\omega_j/\omega_0$,
and
\begin{equation}
 a_j=\sqrt{\frac{E_0}{\omega_0\omega_j}}X_j,
 \qquad
 \pi_j=\sqrt{\frac{E_0\omega_j}{\omega_0}}Y_j.
 \label{eq:unified-scale}
\end{equation}
The product of each position--momentum scale is $E_0/\omega_0$;
division of energy by $E_0$ and rescaling time therefore give the standard
dimensionless canonical equations. Define
\begin{equation}
 S=\Phi_n\operatorname{diag}\sqrt{\frac{E_0}{\omega_0\omega_j}},
 \quad d=S^{\mathsf T}e_\vartheta,\quad
 \ell(s)^{\mathsf T}=N(s)S,
\end{equation}
where $N(s)$ is the Hermite displacement row. Thus the full centreline
and tip slope are recovered in SI units as
$w(s)=\ell(s)^{\mathsf T}X$ and $\vartheta=d^{\mathsf T}X$.
The single dimensionless Hamiltonian used below is
\begin{align}
 h_*(X,Y)&=h_2(X,Y)+R_{\rm mag}(X),\\
 h_2&=\tfrac12\sum_{j=1}^n\nu_j(X_j^2+Y_j^2),\\
 R_{\rm mag}&=1-\cos(d^{\mathsf T}X)-\tfrac12(d^{\mathsf T}X)^2.
 \label{eq:unified-h}
\end{align}
The subtracted quadratic avoids counting the static magnetic stiffness
twice. For the two-mode model,
$(\omega_1,\omega_2)=(11.41728728,70.24049294)\,\mathrm{rad\,s^{-1}}$.
All learned models use this fixed chart; states are not independently
renormalized after fitting.

With $u=B_y/B_0$, $\gamma=\alpha/\omega_0$ and
$\mathsf N=\operatorname{diag}(\nu_j)$, the physical port decoder becomes
\begin{align}
 X'&=\partial_Y h,\\
 Y'&=-\partial_X h-\gamma\mathsf N^{-1}\partial_Yh
          +u\,d\cos(d^{\mathsf T}X).
 \label{eq:unified-port}
\end{align}
For $h=h_*$, the damping term is exactly $-\gamma Y$.
The Rayleigh head is
$\mathcal D(v)=\tfrac12\gamma v^{\mathsf T}\mathsf N^{-1}v$.
For a learned energy, retaining the velocity-based form in
\eqref{eq:unified-port} preserves nonnegative dissipation. Using
$-\gamma Y$ after learning a momentum-dependent residual would, in
general, lose that guarantee. The corresponding power identity is
\begin{equation}
 h'=-\gamma(\partial_Yh)^{\mathsf T}\mathsf N^{-1}\partial_Yh
       +u\cos(d^{\mathsf T}X)d^{\mathsf T}\partial_Yh.
 \label{eq:unified-power}
\end{equation}
Alternatively, including $-u\sin(d^{\mathsf T}X)$ in the time-dependent
energy moves the input power to $-u'\sin(d^{\mathsf T}X)$.
These are equivalent bookkeeping conventions.

\subsection{Independent COMSOL realization of the same discretization}
The model was also constructed and solved in COMSOL Multiphysics
6.3.0.290 through LiveLink for MATLAB. Its two-dimensional Beam interface
uses the Euler--Bernoulli formulation, 32 uniform elements, a consistent
mass matrix, a fixed base, and the nonlinear tip moment
\eqref{eq:unified-tip-torque}. Geometric nonlinearity is disabled to match
the declared small-slope mechanical model. The zero axial displacement
subspace decouples, leaving 64 transverse displacement/rotation degrees
of freedom.

After selecting transverse degrees of freedom by their field names and
coordinates, the independently assembled COMSOL and Python mass and
elastic stiffness matrices have relative Frobenius differences
$7.76\times10^{-15}$ and $3.05\times10^{-14}$, respectively. The first
six modal frequencies agree to a maximum relative difference of
$1.88\times10^{-9}$. These checks establish agreement of the stated
discretizations; they do not establish the material validity of the
chosen parameters. The native model, assembly exports, parameter values
and time-domain verification records accompany the manuscript.

\subsection{A global result for the unforced physical model}
\begin{proposition}[Unforced global stability of the full FE target]
\label{prop:unified-global-stability}
For the full clamped--free Hermite model~\eqref{eq:unified-SI-H}, set
$B_y=0$. With the declared parameters, the potential is globally
strongly convex and the origin is its unique equilibrium.
If $\alpha>0$, the origin is globally asymptotically stable.
If $\alpha=0$, every finite-energy trajectory remains in a bounded
energy sublevel and the origin is Lyapunov stable.
These are global statements for the declared small-slope mathematical
model; they do not establish its physical validity at large deflections
or a global property of arbitrary learned quartic residuals.
\end{proposition}

\begin{proof}
Write $b=e_\vartheta$ and
$U(q)=q^{\mathsf T}Kq/2+E_0\{1-\cos(b^{\mathsf T}q)\}$.
The clamped conditions give
$|w_s(L)|^2\leq L\int_0^L|w_{ss}|^2\,ds$, including for every
Hermite finite-element function. Equivalently,
$bb^{\mathsf T}\preceq [L/(EI)]K$; equality of the corresponding
tip compliance follows from the static quadratic bending shape.
Thus, for every $q$,
\begin{align}
 D^2U(q)&=K+E_0\cos(b^{\mathsf T}q)bb^{\mathsf T}\notag\\
 &\succeq(1-\kappa)K\succeq0.905K,\\
 \kappa&=\frac{E_0L}{EI}=\frac8{27\pi}<0.095.
\end{align}
In particular, $\nabla U=0$ only at $q=0$.
Both $M$ and $K$ are positive definite, and
\[
 H_0\geq\tfrac12p^{\mathsf T}M^{-1}p+\tfrac12q^{\mathsf T}Kq.
\]
Its sublevels are compact. Along the unforced dynamics,
\[
 \dot H_0=-\alpha p^{\mathsf T}M^{-1}p\leq0.
\]
This proves boundedness, global continuation and Lyapunov stability.
For $\alpha>0$, the largest invariant set with $\dot H_0=0$ has
$p=0$ and $\nabla U(q)=0$, hence consists only of the origin.
LaSalle's invariance principle on each compact sublevel gives global
asymptotic stability.
\end{proof}
